% !PRINTED 197
\weberpage{197}
%<<CONT
gesetzt haben) und das obere oder untere Zeichen gilt, je nachdem
$\lambda$ gerade oder ungerade ist. Im letzteren Falle ist dieser
Ausdruck durch $\xi^3 - \eta^3$ theilbar und der Quotient ist durch
$\xi^3 + \eta^3$ und $\xi\eta$ rational und ganzzahlig ausdrückbar (weil eine
symmetrische Function von zwei Grössen so durch die Summe
und das Product dargestellt werden kann). Setzen wir also $\tau = 0$
oder $= 1$, je nachdem $\lambda$ gerade oder ungerade ist, so ergiebt
sich für $C'$ ein Ausdruck von der Form
\[
C'(\xi, \eta) = (\xi^3 - \eta^3)^\tau \; \Sigma \, M (\xi\eta)^\alpha (\xi^3 + \eta^3)^\beta,
\]
worin die $M$ ganzzahlige Coëfficienten sind und $\alpha$, $\beta$ eine Reihe
% [sic] Der Druck sagt "positiver" ganzer Zahlen; als Exponenten können alpha und beta jedoch auch 0 sein, also sind nichtnegative ganze Zahlen gemeint.
positiver ganzer Zahlen durchlaufen, die der Bedingung
\weq{3}{3\tau + 3\beta + 2\alpha = \nu}
genügen. Aus (1) ergiebt sich also, wenn wir mit einer hinlänglich hohen Potenz von 3 multipliciren und $\xi\eta$, $\xi^3+\eta^3$, $\xi^3-\eta^3$
nach \S.\,62, (10) durch $H$, $f$, $Q$ und $r$ durch $D^{-\frac{1}{6}}$ ausdrücken,
\weq{4}{3^\sigma C = Q^\tau \, \Sigma \, M D^\gamma H^\alpha f^\beta,}
worin
\weq{5}{6\gamma = \lambda - 3\tau - 2\alpha,}
und worin die $M$ gleichfalls ganzzahlige Factoren sind.

Zu jedem Werth von $\gamma$ gehört nach (5) ein bestimmter
Werth von $\alpha$ und nach (3) ein bestimmter Werth von $\beta$; ebenso
sind durch einen Werth von $\alpha$ oder von $\beta$ jedesmal die beiden
übrigen Zahlen bestimmt, so dass in (4) jeder Exponent von
$D$, $H$ oder $f$ nur einmal vorkommt.

Dass $\gamma$ eine ganze Zahl ist, folgt leicht aus (2), (3), (5);
denn da nach der Definition $\lambda - 3\tau$ eine gerade Zahl ist, so
ist zunächst $3\gamma$ eine ganze Zahl, und da ferner nach (3)
\[
\lambda - 2\alpha = \frac{3}{2}(\mu - \nu + 2\beta + 2\tau),
\]
so ist $6\gamma$ und mithin $3\gamma$ durch 3 theilbar. Hieraus folgt, dass
% [sic] Der Druck nennt hier sigma. Die unmittelbar vorhergehende Teilbarkeitsrechnung beweist jedoch, dass gamma ganzzahlig ist; sigma war bereits als genügend hoher Exponent von 3 eingeführt.
auch $\sigma$ eine ganze Zahl ist.

Dass aber $\gamma$ auch positiv sein muss, sehen wir so ein:

Angenommen, auf der rechten Seite von (4) kommen auch
negative Potenzen von $D$ vor, so multipliciren wir diese Gleichung
beiderseits mit einer so hohen Potenz von $D$, dass rechts keine
negativen Potenzen von $D$ mehr auftreten, dass aber auch die
rechte Seite nicht den Factor $D$ erhält. Setzen wir in der so
gewonnenen Gleichung
%CONT>>
